% questions.tex —— 2011 年·天津卷·文科 · 题目内容片段（无答案）
% 由 tools/generate_projects.py 依来源总账自 2011/2011_数学真题_天津文_试题.md 生成，勿手改（改源 MD 后重跑）

\gkpart{一、选择题}

\begin{question}
\gkti{1}$\mathrm i$ 是虚数单位，复数 $\dfrac{1-3\mathrm i}{1-\mathrm i}=$\mbox{（　　）}

（A）$2-\mathrm i$；（B）$2+\mathrm i$；（C）$-1-2\mathrm i$；（D）$-1+2\mathrm i$。
\end{question}
\begin{question}
\gkti{2}设变量 $x$，$y$ 满足约束条件 $\begin{cases}x\geqslant1,\\ x+y-4\leqslant0,\\ x-3y+4\leqslant0,\end{cases}$ 则目标函数 $z=3x-y$ 的最大值为\mbox{（　　）}

（A）$-4$；（B）$0$；（C）$\dfrac43$；（D）$4$。
\Needspace{0.4\textheight}
\end{question}
\begin{question}
\gkti{3}阅读程序框图，运行相应的程序，若输入 $x$ 的值为 $-4$，则输出 $y$ 的值为\mbox{（　　）}
\nopagebreak[4]
\begin{center}% 2011 年天津卷文科第 3 题：程序框图（循环：|x|>3 时执行 x=|x-3| 再回判断）
% 来源 PDF：普通高考/2011/2011天津文.pdf，第 1 页（印面 707）第 1 栏。
% 重绘依据：原图为「开始 → 输入 x → 判断 |x|>3?」，是分支向下执行 x=|x-3| 并回到判断，
%   否分支向右再向下汇入 y=2^x，随后「输出 y」与「结束」。
\begin{tikzpicture}[scale=0.72, semithick, >=Stealth,
  block/.style={draw, minimum width=3.2cm, minimum height=0.6cm, align=center},
  dec/.style={draw, diamond, aspect=3.0, inner sep=1pt, align=center},
  arr/.style={->}]
  \node (start) [draw, rounded corners=6pt, minimum width=1.5cm, minimum height=0.5cm] {开始};
  \node (in)    [block, below=0.35cm of start] {输入 $x$};
  \node (dec)   [dec, below=0.5cm of in] {$|x|>3$?};
  \node (op)    [block, below=0.55cm of dec] {$x=|x-3|$};
  \node (pow)   [block, below=1.0cm of op] {$y=2^x$};
  \node (out)   [draw, trapezium, trapezium left angle=70, trapezium right angle=110,
                 minimum width=3.0cm, minimum height=0.6cm, below=0.45cm of pow] {输出 $y$};
  \node (end)   [draw, rounded corners=6pt, minimum width=1.5cm, minimum height=0.5cm,
                 below=0.4cm of out] {结束};
  \draw[arr] (start) -- (in);
  \draw[arr] (in) -- (dec);
  \draw[arr] (dec) -- node[right] {是} (op);
  \draw[arr] (op.west) -- ++(-1.6,0) |- (dec.west);
  \draw[arr] (dec.east) -- node[above] {否} ++(1.9,0) |- (pow.east);
  \draw[arr] (op) -- (pow);
  \draw[arr] (pow) -- (out);
  \draw[arr] (out) -- (end);
\end{tikzpicture}
\end{center}
\nopagebreak[4]
（A）$0.5$；（B）$1$；（C）$2$；（D）$4$。
\end{question}
\begin{question}
\gkti{4}设集合 $A=\{x\in\mathbf R\mid x-2>0\}$，$B=\{x\in\mathbf R\mid x<0\}$，$C=\{x\in\mathbf R\mid x(x-2)>0\}$，则「$x\in A\cup B$」是「$x\in C$」的\mbox{（　　）}

（A）充分而不必要条件；（B）必要而不充分条件；（C）充分必要条件；（D）既不充分也不必要条件。
\end{question}
\begin{question}
\gkti{5}已知 $a=\log_23.6$，$b=\log_43.2$，$c=\log_43.6$，则\mbox{（　　）}

（A）$a>b>c$；（B）$a>c>b$；（C）$b>a>c$；（D）$c>a>b$。
\end{question}
\begin{question}
\gkti{6}已知双曲线 $\dfrac{x^2}{a^2}-\dfrac{y^2}{b^2}=1$（$a>0$，$b>0$）的左顶点与抛物线 $y^2=2px$（$p>0$）的焦点的距离为 $4$，且双曲线的一条渐近线与抛物线的准线的交点坐标为 $(-2,-1)$，则双曲线的焦距为\mbox{（　　）}

（A）$2\sqrt3$；（B）$2\sqrt5$；（C）$4\sqrt3$；（D）$4\sqrt5$。
\end{question}
\begin{question}
\gkti{7}已知函数 $f(x)=2\sin(\omega x+\varphi)$，$x\in\mathbf R$，其中 $\omega>0$，$-\pi<\varphi\leqslant\pi$。若 $f(x)$ 的最小正周期为 $6\pi$，且当 $x=\dfrac\pi2$ 时，$f(x)$ 取得最大值，则\mbox{（　　）}

（A）$f(x)$ 在区间 $[-2\pi,0]$ 上是增函数；（B）$f(x)$ 在区间 $[-3\pi,-\pi]$ 上是增函数；（C）$f(x)$ 在区间 $[3\pi,5\pi]$ 上是减函数；（D）$f(x)$ 在区间 $[4\pi,6\pi]$ 上是减函数。
\end{question}
\begin{question}
\gkti{8}对实数 $a$ 和 $b$，定义运算「$\otimes$」：$a\otimes b=\begin{cases}a,&a-b\leqslant1,\\ b,&a-b>1.\end{cases}$ 设函数 $f(x)=(x^2-2)\otimes(x-1)$，$x\in\mathbf R$。若函数 $y=f(x)-c$ 的图象与 $x$ 轴恰有两个公共点，则实数 $c$ 的取值范围是\mbox{（　　）}

（A）$(-1,1]\cup(2,+\infty)$；（B）$(-2,-1]\cup(1,2]$；（C）$(-\infty,-2)\cup(1,2]$；（D）$[-2,-1]$。
\end{question}
\gkpart{二、填空题}

\begin{question}
\gkti{9}已知集合 $A=\{x\in\mathbf R\mid |x-1|<2\}$，$\mathbf Z$ 为整数集，则集合 $A\cap\mathbf Z$ 中所有元素的和等于\gkblank{}。

\Needspace{0.4\textheight}
\end{question}
\begin{question}
\gkti{10}一个几何体的三视图如图所示（单位：$\mathrm m$），则这个几何体的体积为\gkblank{}$\mathrm m^3$。
\nopagebreak[4]
\begin{center}% 2011 年天津卷文科第 10 题：三视图（单位：m）
% 来源 PDF：普通高考/2011/2011天津文.pdf，第 1 页（印面 707）第 2 栏。
% 重绘依据：正视图总宽 2、总高 3，上块宽 1 高 2，下块宽 2 高 1；
%   侧视图宽 1、高 3，在距底边 1 处有一条水平分界线；俯视图宽 2、深 1，在 x=1 处有分界线。
\begin{tikzpicture}[scale=0.78, semithick, every node/.style={font=\small}]
  % ---------- 正视图 ----------
  \begin{scope}[shift={(0,0)}]
    \draw (0,0) rectangle (2,1);          % 下块 2×1
    \draw (0,1) rectangle (1,3);          % 上块 1×2
    \draw[<->] (-0.35,0) -- (-0.35,3) node[midway,left] {$2$};
    \draw[<->] (0,-0.35) -- (2,-0.35) node[midway,below] {$2$};
    \draw[<->] (1.35,0) -- (1.35,1) node[midway,right] {$1$};
    \node at (1,-1.0) {正视图};
  \end{scope}
  % ---------- 侧视图 ----------
  \begin{scope}[shift={(4.4,0)}]
    \draw (0,0) rectangle (1,3);
    \draw (0,1) -- (1,1);
    \draw[<->] (-0.35,0) -- (-0.35,3) node[midway,left] {$2$};
    \draw[<->] (1.3,0) -- (1.3,1) node[midway,right] {$1$};
    \draw[<->] (0,-0.35) -- (1,-0.35) node[midway,below] {$1$};
    \node at (0.5,-1.0) {侧视图};
  \end{scope}
  % ---------- 俯视图 ----------
  \begin{scope}[shift={(0,-4.6)}]
    \draw (0,0) rectangle (2,1);
    \draw (1,0) -- (1,1);
    \draw[<->] (0,-0.35) -- (2,-0.35) node[midway,below] {$2$};
    \draw[<->] (2.3,0) -- (2.3,1) node[midway,right] {$1$};
    \draw[<->] (0,1.35) -- (1,1.35) node[midway,above] {$1$};
    \node at (1,-1.0) {俯视图};
  \end{scope}
\end{tikzpicture}
\end{center}
\nopagebreak[4]
\end{question}
\begin{question}
\gkti{11}已知 $\{a_n\}$ 为等差数列，$S_n$ 为其前 $n$ 项和，$n\in\mathbf N^*$。若 $a_3=16$，$S_{20}=20$，则 $S_{10}$ 的值为\gkblank{}。

\end{question}
\begin{question}
\gkti{12}已知 $\log_2a+\log_2b\geqslant1$，则 $3^a+9^b$ 的最小值为\gkblank{}。

\Needspace{0.4\textheight}
\end{question}
\begin{question}
\gkti{13}如图，已知圆中两条弦 $AB$ 与 $CD$ 相交于点 $F$，$E$ 是 $AB$ 延长线上一点，且 $DF=CF=\sqrt2$，$AF:FB:BE=4:2:1$。若 $CE$ 与圆相切，则 $CE$ 的长为\gkblank{}。
\nopagebreak[4]
\begin{center}% 2011 年天津卷文科第 13 题：圆中两弦 AB 与 CD 相交于 F，E 在 AB 延长线上，CE 与圆相切。
% 来源 PDF：普通高考/2011/2011天津文.pdf，第 1 页（印面 707）第 2 栏。
% 重绘依据：原图圆心约在 (0,0)，A、B 在圆上且 AB 近水平，E 在 B 的右侧圆外，
%   D 在圆上左上方，C 在圆上右下方，CD 过 F；CE 为切线，C 为切点。
\begin{tikzpicture}[scale=1.0, semithick, every node/.style={inner sep=1.5pt}]
  \coordinate (O) at (0,0);
  \draw (O) circle (2.0);
  % 弦 AB（近水平，略向右上倾斜）
  \coordinate (A) at (-160:2.0); % 左下
  \coordinate (B) at (-15:2.0);  % 右下
  \draw (A) -- (B);
  % E 在 AB 延长线上
  \coordinate (E) at (2.55,-0.28);
  \draw (B) -- (E);
  % 弦 CD 过 F（F 为 AB 与 CD 交点）
  \coordinate (F) at ($(A)!0.62!(B)$);
  \coordinate (D) at (135:2.0);
  \coordinate (C) at (-58:2.0);
  \draw (D) -- (C);
  % 切线 CE
  \draw (C) -- (E);
  % 标注
  \node[left]  at (A) {$A$};
  \node[above right] at (B) {$B$};
  \node[right] at (E) {$E$};
  \node[above left] at (D) {$D$};
  \node[below right] at (C) {$C$};
  \node[above] at (F) {$F$};
\end{tikzpicture}
\end{center}
\nopagebreak[4]
\end{question}
\begin{question}
\gkti{14}已知直角梯形 $ABCD$ 中，$AD\parallel BC$，$\angle ADC=90^\circ$，$AD=2$，$BC=1$，$P$ 是腰 $DC$ 上的动点，则 $|\overrightarrow{PA}+3\overrightarrow{PB}|$ 的最小值为\gkblank{}。

\end{question}
\gkpart{三、解答题}

\begin{question}
\gkti{15}编号分别为 $A_1$，$A_2$，$\cdots$，$A_{16}$ 的 $16$ 名篮球运动员在某次训练比赛中的得分记录如下：

\begin{center}
\begin{tabular}{lllll}
\toprule
编号 & $A_1$ & $A_2$ & $A_3$ & $A_4$ \\
\midrule
得分 & $15$ & $35$ & $21$ & $28$ \\
\bottomrule
\end{tabular}
\end{center}

\begin{center}
\begin{tabular}{lllll}
\toprule
编号 & $A_5$ & $A_6$ & $A_7$ & $A_8$ \\
\midrule
得分 & $25$ & $36$ & $18$ & $34$ \\
\bottomrule
\end{tabular}
\end{center}

\begin{center}
\begin{tabular}{lllll}
\toprule
编号 & $A_9$ & $A_{10}$ & $A_{11}$ & $A_{12}$ \\
\midrule
得分 & $17$ & $26$ & $25$ & $33$ \\
\bottomrule
\end{tabular}
\end{center}

\begin{center}
\begin{tabular}{lllll}
\toprule
编号 & $A_{13}$ & $A_{14}$ & $A_{15}$ & $A_{16}$ \\
\midrule
得分 & $22$ & $12$ & $31$ & $38$ \\
\bottomrule
\end{tabular}
\end{center}

（1） 将得分在对应区间内的人数填入相应的空格：

\begin{center}
\begin{tabular}{llll}
\toprule
区间 & [10,20) & [20,30) & [30,40] \\
\midrule
人数 &  &  &  \\
\bottomrule
\end{tabular}
\end{center}

（2） 从得分在区间 $[20,30)$ 内的运动员中随机抽取 $2$ 人，
① 用运动员编号列出所有可能的抽取结果；
② 求这 $2$ 人得分之和大于 $50$ 的概率。

\end{question}
\begin{question}
\gkti{16}在 $\triangle ABC$ 中，内角 $A$，$B$，$C$ 的对边分别为 $a$，$b$，$c$，已知 $B=C$，$2b=\sqrt3a$。

（1） 求 $\cos A$ 的值；

（2） 求 $\cos\left(2A+\dfrac\pi4\right)$ 的值。

\Needspace{0.4\textheight}
\end{question}
\begin{question}
\gkti{17}如图，在四棱锥 $P-ABCD$ 中，底面 $ABCD$ 为平行四边形，$\angle ADC=45^\circ$，$AD=AC=1$，$O$ 为 $AC$ 中点，$PO\perp$ 平面 $ABCD$，$PO=2$，$M$ 为 $PD$ 中点。

（1） 证明：$PB\parallel$ 平面 $ACM$；

（2） 证明：$AD\perp$ 平面 $PAC$；

（3） 求直线 $AM$ 与平面 $ABCD$ 所成角的正切值。
\nopagebreak[4]
\begin{center}% 2011 年天津卷文科第 17 题：四棱锥 P-ABCD，底面 ABCD 为平行四边形，
% O 为 AC 中点，PO⊥底面，M 为 PD 中点。
% 来源 PDF：普通高考/2011/2011天津文.pdf，第 2 页（印面 708）第 1 栏。
% 重绘依据：原图为斜二测式透视，D 在左后、A 在左前、B 在右前、C 在右后；
%   P 在 O 正上方；可见棱 PA、PB、PC、PD、AB、BC 画实线，被遮挡的 AD、CD、AC
%   与辅助线 OM 画虚线。
\begin{tikzpicture}[scale=0.95, semithick,
  every node/.style={inner sep=1.5pt}]
  \coordinate (A) at (0.30,0.20);
  \coordinate (B) at (3.30,0.20);
  \coordinate (C) at (4.10,1.35);
  \coordinate (D) at (1.10,1.35);
  \coordinate (O) at (2.20,0.78);
  \coordinate (P) at (2.20,3.70);
  \coordinate (M) at ($(P)!0.5!(D)$);
  % 底面被遮挡的边与对角线
  \draw[dashed] (A) -- (D) -- (C);
  \draw[dashed] (A) -- (C);
  % 可见棱
  \draw (A) -- (B) -- (C);
  \draw (P) -- (A);
  \draw (P) -- (B);
  \draw (P) -- (C);
  \draw (P) -- (D);
  % 辅助线 OM（中位线，图中虚线）
  \draw[dashed] (O) -- (M);
  % 标注
  \node[left] at ([xshift=-1pt]A) {$A$};
  \node[below] at (B) {$B$};
  \node[right] at (C) {$C$};
  \node[above left] at (D) {$D$};
  \node[above] at (P) {$P$};
  \node[left] at ([xshift=1pt]M) {$M$};
  \node[below] at (O) {$O$};
\end{tikzpicture}
\end{center}
\nopagebreak[4]
\end{question}
\begin{question}
\gkti{18}设椭圆 $\dfrac{x^2}{a^2}+\dfrac{y^2}{b^2}=1$（$a>b>0$）的左、右焦点分别为 $F_1$，$F_2$。点 $P(a,b)$ 满足 $|PF_2|=|F_1F_2|$。

（1） 求椭圆的离心率 $e$；

（2） 设直线 $PF_2$ 与椭圆相交于 $A$，$B$ 两点，若直线 $PF_2$ 与圆 $(x+1)^2+(y-\sqrt3)^2=16$ 相交于 $M$，$N$ 两点，且 $|MN|=\dfrac58|AB|$，求椭圆的方程。

\end{question}
\begin{question}
\gkti{19}已知函数 $f(x)=4x^3+3tx^2-6t^2x+t-1$，$x\in\mathbf R$，其中 $t\in\mathbf R$。

（1） 当 $t=1$ 时，求曲线 $y=f(x)$ 在点 $(0,f(0))$ 处的切线方程；

（2） 当 $t\neq0$ 时，求 $f(x)$ 的单调区间；

（3） 证明：对任意的 $t\in(0,+\infty)$，$f(x)$ 在区间 $(0,1)$ 内均存在零点。

\end{question}
\begin{question}
\gkti{20}已知数列 $\{a_n\}$ 与 $\{b_n\}$ 满足 $b_{n+1}a_n+b_na_{n+1}=(-2)^n+1$，$b_n=\dfrac{3+(-1)^{n-1}}2$，$n\in\mathbf N^*$，且 $a_1=2$。

（1） 求 $a_2$，$a_3$ 的值；

（2） 设 $c_n=a_{2n+1}-a_{2n-1}$，$n\in\mathbf N^*$，证明：数列 $\{c_n\}$ 是等比数列；

（3） 设 $S_n$ 为 $\{a_n\}$ 的前 $n$ 项和，证明：$\dfrac{S_1}{a_1}+\dfrac{S_2}{a_2}+\cdots+\dfrac{S_{2n-1}}{a_{2n-1}}+\dfrac{S_{2n}}{a_{2n}}\leqslant n-\dfrac13$（$n\in\mathbf N^*$）。

\clearpage

\end{question}